% Partie : logique intuitionniste
% Chapitre : propositions comme types
% Section : introduction

\documentclass[../../../include/open-logic-section]{subfiles}

\begin{document}

\olfileid[fr]{int}{pty}{int}

\olsection{Introduction}

Le lambda-calcul non typé est un modèle de calcul simple. Il manipule
des termes construits à partir de variables. Si $N$ et $M$
sont des termes, alors $(NM)$ est aussi un terme :
« $N$ appliqué à~$M$ ». Si $N$ est un terme,
$\lambd[x][N]$ en est un également. Les occurrences libres
de la variable~$x$ dans $N$ deviennent liées par l'opérateur
$\lambd[x]$ dans $\lambd[x][N]$.
Un terme de la forme $(\lambd[x][N]M)$ se
$\beta$-contracte en $\Subst{N}{M}{x}$ : on remplace
chaque occurrence libre de $x$ dans~$N$ par $M$,
en renommant les variables liées si nécessaire pour éviter
la capture. Un terme $N$ se réduit en une étape en~$N'$
si l'on obtient $N'$ en $\beta$-contractant un sous-terme
de $N$ ; on écrit $N \redone N'$.
Un calcul dans le $\lambd$-calcul est une suite
$N \redone N_1 \redone N_2 \redone \dots$.
Si cette suite aboutit à un terme $N_k$ dont aucun sous-terme
ne peut être $\beta$-contracté, ce terme est dit \emph{normal},
ou constitue une forme normale de~$N$.
On représente ainsi un calcul par une suite de réductions.
Il s'arrête lorsqu'il atteint une forme normale, qui en donne
le résultat. Ce modèle simple est Turing-complet : toute
fonction partielle récursive peut être calculée par un terme
du $\lambd$-calcul.

Haskell Curry et William Howard ont découvert indépendamment
une correspondance étroite entre la déduction naturelle et le
$\lambd$-calcul. Intuitivement, un terme $\lambd[x][N]$
est une fonction : $x$ représente l'argument, et $N$ indique
comment en calculer la valeur à partir de~$x$.
Le terme $(\lambd[x][N]M)$ applique donc la fonction
$\lambd[x][N]$ à un argument~$M$. La $\beta$-contraction
indique comment évaluer cette application : remplacer
$x$ par~$M$ dans $N$, puis poursuivre l'évaluation du résultat.
Considérons maintenant des fonctions dont le domaine et le
codomaine sont fixés. Si $\lambd[x][N]$ a pour domaine $A$
et pour codomaine~$B$, la variable~$x$ doit recevoir
des arguments de~$A$, et $N$ doit produire des valeurs de~$B$.
L'expression $(\lambd[x][N]M)$ ne convient donc que si
$\lambd[x][N]$ est une fonction $A \to B$ et $M \in A$.
Ajouter cette information au $\lambd$-calcul conduit au
$\lambd$-calcul \emph{typé}.
Toute expression $(\lambd[x][N]M)$ n'y est pas admise :
les types de $\lambd[x][N]$ et de $M$ doivent correspondre.
Le premier terme doit avoir le type $A \lif B$, le second
le type~$A$. Le terme $\lambd[x][N]$ a le type $A \to B$
si $x$ a le type $A$ et $N$ le type~$B$.
Plus généralement, une expression $(M_1M_2)$ est admise
lorsque $M_1$ a le type $A \to B$ et $M_2$ le type~$A$ ;
$(M_1M_2)$ a alors le type~$B$.

Ces règles de typage correspondent aux dérivations en
déduction naturelle : les types y sont représentés par des
formules, et les dérivations par des $\lambd$-termes.
Par exemple, une dérivation de $!A \lif !B$ correspond
à un terme $\lambd[\typeof{x}{!A}][N]$ de type fonction
$!A \lif !B$.
Les règles d'inférence de la déduction naturelle correspondent
aux règles de typage du lambda-calcul typé. Ainsi :
\begin{prooftree}
  \AxiomC{$\Discharge{!A}{x}$}
  \DeduceC{$!B$}
  \DischargeRule{\Intro\lif}{x}
  \UnaryInfC{$!A \lif !B$}
  \DisplayProof\bottomAlignProof
  \quad\text{correspond à}\quad
  \Axiom$x: !A \fCenter N: !B$
  \UnaryInf$ \fCenter \lambd[\typeof{x}{!A}][N] : !A \lif !B$
\end{prooftree}
Si $N$ a le type~$!B$ sous l'hypothèse que $x$ a le
type~$!A$, alors $\lambd[\typeof{x}{!A}][N]$
a le type~$!A \lif !B$.

La règle $\Elim\lif$ correspond à la règle de typage
des termes d'application :
\[
  \AxiomC{$!A \lif !B$}
  \AxiomC{$!A$}
  \RightLabel{\Elim\lif}
  \BinaryInfC{$!B$}
  \DisplayProof
  \quad\text{correspond à}\quad
  \Axiom$\fCenter M_1: !A \lif !B$
  \Axiom$\fCenter M_2: !A$
  \BinaryInf$ \fCenter (M_1M_2) : !B$
  \DisplayProof
\]
Lorsqu'une règle \Intro\lif{} est immédiatement suivie
d'une règle \Elim\lif{}, la dérivation peut être simplifiée :
\begin{gather*}
  \AxiomC{$\Discharge{!A}{x}$}
  \DeduceC{$!B$}
  \DischargeRule{\Intro\lif}{x}
  \UnaryInfC{$!A \lif !B$}
  \AxiomC{}
  \DeduceC{$!A$}
  \RightLabel{\Elim\lif}
  \BinaryInfC{$!B$}
  \DisplayProof
  \quad\redone\quad
  \AxiomC{}
  \DeduceC{$!A$}
  \DeduceC{$!B$}
  \DisplayProof
\end{gather*}
Cette simplification correspond à la $\beta$-contraction
des termes lambda :
\[
(\lambd[\typeof{x}{!A}][N]M) \quad\redone\quad
\Subst{N}{M}{x}.
\]

Des correspondances analogues relient les règles pour
$\land$ aux types « produit », et les règles pour $\lor$
aux types « somme ».

La correspondance entre les termes du lambda-calcul simplement
typé et les dérivations en déduction naturelle porte les noms
de « correspondance de Curry--Howard », « preuves comme
programmes » ou « propositions comme types ».
Outre les formules, ou propositions, qui correspondent aux
types, et les preuves qui correspondent aux termes,
on peut résumer ces correspondances ainsi :
\begin{center}
  \begin{tabular}{c c c}
    logique & programme \\
    \hline
    proposition/formule & type \\
    preuve & terme \\
    hypothèse & variable \\
    hypothèse déchargée & variable liée \\
    hypothèse non déchargée & variable libre \\
    $!A \lif !B$ & type fonction \\
    $!A \land !B$ & type produit \\
    $!A \lor !B$ & type somme \\
    $\lfalse$ & type vide \\
    \hline
  \end{tabular}
\end{center}

La correspondance de Curry--Howard est l'un des fondements
des assistants de preuve et des vérificateurs de types :
vérifier une preuve qui témoigne d'une proposition revient
à vérifier qu'un programme, représenté par un terme,
possède le type annoncé.

\end{document}
